\section{Conclusion}

\subsection{Summary}

We presented \textbf{DecayLinearAttention}, a linear attention module with 
\textbf{data-dependent exponential decay per head} and **chunked scan with 
O(1) state memory**. The method supports bidirectional and causal modes, 
GQA, RoPE, and cross-attention.

\subsection{Results Recap}

\begin{table}[H]
\centering
\small
\begin{tabular}{llll}
\toprule
\textbf{Model} & \textbf{Acc $\Delta$} & \textbf{State mem} & \textbf{Time ratio at L=2048} \\
\midrule
DistilBERT & {{D\_DELTA}}\% & \textbf{16 KB} & {{D\_RATIO\_2048}}$\times$ \\
RoBERTa & {{R\_DELTA}}\% & \textbf{16 KB} & {{R\_RATIO\_2048}}$\times$ \\
GPT-2 & {{G\_DELTA}}\% & \textbf{16 KB} & {{G\_RATIO\_2048}}$\times$ \\
\bottomrule
\end{tabular}
\end{table}

\subsection{Key Takeaways}

\begin{enumerate}
\item \textbf{Truly O(1) state memory.}
\item \textbf{Data-dependent decay improves accuracy.}
\item \textbf{Per-head bias gives multi-scale horizons.}
\item \textbf{Cross-attention works.}
\item \textbf{Speedup grows with L.}

\end{enumerate}
\subsection{Future Work}

\subsubsection{Short-Term}
\begin{itemize}
\item Fused Triton kernels for chunked scan
\item Adaptive chunk size
\item Learned feature map

\end{itemize}
\subsubsection{Long-Term}
\begin{itemize}
\item Scale to L > 32K
\item Combine with token compression
\item Theoretical analysis

\end{itemize}
\subsection{Closing Remarks}

DecayLinearAttention offers a practical, memory-efficient alternative to 
softmax attention for long sequences. We release all code and checkpoints.

\subsection{Acknowledgments}

We thank {{ACKNOWLEDGMENTS}}.

\subsection{Author Contributions}

\begin{itemize}
\item \textbf{{{AUTHOR\_1}}} designed the module, implemented, ran experiments, wrote
  the paper.
\item \textbf{{{AUTHOR\_2}}} provided feedback, edited the manuscript.

\end{itemize}
\subsection{Conflict of Interest}

None.
\begin{lstlisting}

---
\end{lstlisting}